% fig_arch.tex
% Architecture diagram. Compact horizontal layout:
% (LEFT)  conditioning chain     -> FiLM input to transformer
% (CENTRE) data flow column      -> shifted input -> transformer -> head
% (RIGHT) inference render chain -> sample -> renderer -> output

\begin{figure}[t]
\centering
\resizebox{\columnwidth}{!}{%
\begin{tikzpicture}[
    font=\sffamily\footnotesize,
    >=Stealth,
    node distance=3mm,
    box/.style    ={draw, rounded corners=2pt, inner sep=3pt,
                    minimum height=7mm, minimum width=22mm,
                    align=center, fill=white},
    txtbox/.style ={box, fill=blue!8},
    netbox/.style ={box, fill=orange!12},
    headbox/.style={box, fill=green!12},
    rendbox/.style={box, fill=red!10},
    inbox/.style  ={box, fill=gray!10},
    arr/.style={->, semithick},
]

% =========================================================
% Centre column: data flow (top -> bottom)
% =========================================================
\node[inbox]              (sig)   {target $\mathbf{s}\in[0,1]^{N\times C}$};
\node[inbox, below=of sig](input) {teacher-forced input};
\node[box,   fill=yellow!15, below=of input] (proj)
        {input projection};
\node[box,   fill=yellow!15, below=of proj] (pe)
        {{\bfseries +} 2D path PE};

% Transformer (visual stack)
\begin{scope}[on background layer]
\node[netbox, below=of pe, xshift=0.7mm, yshift=-0.7mm] (b2) {};
\end{scope}
\node[netbox, below=of pe] (tx)
    {causal transformer\\\scriptsize 6 blocks, FiLM};

\node[headbox, below=of tx] (head)
    {DMoL head\\\scriptsize $K{=}5$ mixtures};

% Centre arrows
\draw[arr] (sig.south)   -- (input.north);
\draw[arr] (input.south) -- (proj.north);
\draw[arr] (proj.south)  -- (pe.north);
\draw[arr] (pe.south)    -- (tx.north);
\draw[arr] (tx.south)    -- (head.north);

% =========================================================
% Left column: conditioning, aligned with transformer
% =========================================================
\node[txtbox, left=10mm of pe.west,  anchor=east] (txt)  {prompt};
\node[txtbox, left=10mm of tx.west,  anchor=east] (clip) {CLIP ViT-B/32};
\node[txtbox, left=10mm of head.west,anchor=east] (cond) {cond.\ vector $\mathbf{c}$};
\draw[arr] (txt.south)  -- (clip.north);
\draw[arr] (clip.south) -- (cond.north);

% Conditioning -> transformer (FiLM)
\draw[arr] (cond.east) .. controls +(15mm,0) and +(-12mm,0) ..
    node[pos=0.65, above, font=\sffamily\scriptsize]{FiLM}
    (tx.west);

% =========================================================
% Right column: inference rendering chain
% =========================================================
\node[inbox,   right=12mm of head, anchor=west] (pred)
    {predicted $\hat{\mathbf{s}}$};
\node[rendbox, above=of pred] (rend)
    {Gaussian-beam\\renderer};
\node[inbox, above=of rend] (out)
    {output $\hat{\mathbf{x}}$};

\draw[arr] (head.east) -- node[above, font=\sffamily\scriptsize]{sample}
    (pred.west);
\draw[arr] (pred.north) -- (rend.south);
\draw[arr] (rend.north) -- (out.south);

% =========================================================
% Training loss (dashed): head -> target signal
% =========================================================
\draw[arr, dashed, gray!75!black, thick]
    (head.north east) .. controls +(8mm,4mm) and +(8mm,-4mm) ..
    node[pos=0.55, right, font=\sffamily\scriptsize\itshape, black]
        {NLL loss}
    (sig.east);

\end{tikzpicture}%
}
\caption{Signal-space autoregressive architecture. The target image
is converted to a 1D signal $\mathbf{s}$ by sampling along a fixed
raster path (omitted from the figure for clarity). A causal
transformer with FiLM conditioning on a CLIP text embedding
$\mathbf{c}$ predicts the next signal value at each step. The DMoL
output head provides the training likelihood (dashed arrow back to
the target). At inference, the predicted signal $\hat{\mathbf{s}}$
is mapped to an image $\hat{\mathbf{x}}$ by the differentiable
Gaussian-beam renderer (right column).}
\label{fig:arch}
\end{figure}
